\chapter{Implicaciones de política hospitalaria}
\label{ch:politica}

Los resultados teóricos y empíricos de esta tesis tienen implicaciones directas para el diseño de la política hospitalaria en España y, de forma más general, para cualquier sistema sanitario que contemple la coexistencia de formas organizativas centralizadas y descentralizadas.

Las recomendaciones que se formulan en este capítulo se derivan de las correlaciones robustas identificadas en los capítulos empíricos y de las predicciones del modelo teórico. Deben entenderse como hipótesis directrices para el diseño de políticas hospitalarias, no como prescripciones causales definitivas. El diseño observacional de esta tesis ---un panel de datos con controles regionales y de complejidad, sin asignación aleatoria de la estructura organizativa--- identifica asociaciones condicionales robustas, pero no establece causalidad en sentido estricto. La validación causal de estas recomendaciones requeriría diseños cuasiexperimentales que exploten variación exógena en la asignación de hospitales a estructuras organizativas, como las reversiones de concesiones administrativas en la Comunidad Valenciana \citep[véase][para una discusión de estos diseños]{angrist2009}. No obstante, las correlaciones identificadas, su coherencia con la teoría y su robustez frente a múltiples especificaciones justifican su utilización como base informativa para el diseño institucional.

\section{Sobre la estructura organizativa}

\subsection{La descentralización no es unidimensionalmente superior}

El resultado empírico más robusto de esta tesis es que la descentralización reduce la ineficiencia técnica en la dimensión de cantidad (P1, confirmada en 7/7 especificaciones). Los hospitales con mayor autonomía en la gestión de contratos e incentivos producen más actividad por unidad de input. Sin embargo, esta ventaja en volumen no se extiende de forma simétrica a la dimensión de intensidad, donde los efectos son menos pronunciados y, en el mejor de los casos, asimétricos (P2).

Esta asimetría tiene una consecuencia política directa: \textit{la evaluación de la eficiencia hospitalaria no puede basarse exclusivamente en indicadores de volumen}. Un hospital descentralizado que muestre altas puntuaciones de eficiencia en el número de altas atendidas puede no ser igualmente eficiente en la intensidad diagnóstico-terapéutica por episodio. Los reguladores sanitarios deberían incorporar medidas de intensidad y calidad en los mecanismos de evaluación y pago a los hospitales.

\subsection{La cautela ante las reversiones de concesiones}

En el período analizado (2010--2023), varias comunidades autónomas han revertido concesiones administrativas que habían descentralizado la gestión de hospitales públicos \citep{repullo2008}. Los resultados de esta tesis sugieren que estas reversiones pueden tener costes de eficiencia en la dimensión de cantidad que no están siendo evaluados adecuadamente. Al mismo tiempo, el modelo teórico muestra que la estructura centralizada tiene ventajas en la provisión de intensidad, especialmente cuando las tareas del médico son complementarias (Resultado~\ref{res:ruido}).

La recomendación derivada es que las decisiones sobre la forma organizativa de los hospitales públicos deberían basarse en evaluaciones empíricas de eficiencia multidimensional, no en posiciones ideológicas sobre la titularidad pública o privada de la gestión.

\section{Sobre los mecanismos de pago}

\subsection{La potencia de los incentivos debe calibrarse al output relevante}

El modelo teórico muestra que, bajo sustitución entre tareas, los incentivos de volumen desplazan esfuerzo desde la intensidad hacia la cantidad (Resultado~\ref{res:psi12}). Los datos empíricos confirman esta predicción: la interacción entre estructura descentralizada y tipo de pago ($D^{desc}\times pct^{SNS}$) tiene efectos significativamente distintos sobre cantidad e intensidad (P3 confirmada).

Esto sugiere que los esquemas de pago por actividad (GRD, pago capitativo ajustado) deberían complementarse con indicadores de intensidad y calidad vinculados a parte de la retribución. Los \textit{contratos-programa} autonómicos que incluyen umbrales de intensidad diagnóstica o indicadores de calidad asistencial van en esta dirección, pero el modelo teórico indica que su efectividad depende críticamente de la precisión con la que se miden las distintas dimensiones del output ($\sigma_1$ y $\sigma_2$ en el modelo).

\subsection{Los incentivos al médico deben diseñarse en función de la estructura organizativa}

Una de las conclusiones más interesantes del modelo teórico es que, bajo sustitución perfecta entre tareas, la estructura centralizada óptima no incluye incentivos para el médico (salario fijo), mientras que la descentralizada sí (Resultado~\ref{res:incentivos_centralizados}). Este resultado proporciona un fundamento teórico para la observación empírica de que los hospitales privados utilizan con mayor frecuencia esquemas de retribución variable para los médicos, mientras que el sistema público mantiene predominantemente remuneraciones salariales fijas con complementos débiles.

La implicación política no es que el sistema público deba imitar los incentivos del sector privado, sino que el tipo de incentivo debe ser coherente con la estructura organizativa. Si se mantiene la contratación centralizada del personal médico (como es la norma en el SNS español), la potencia de los incentivos directos al médico debería ser baja, y el financiador debería centrarse en los incentivos al hospital. Si, por el contrario, se descentraliza la gestión, los incentivos al médico adquieren un papel relevante como mecanismo de transmisión.

\section{Sobre la medición del output hospitalario}

\subsection{La necesidad de medir la intensidad}

Los resultados empíricos revelan que la frontera de intensidad tiene niveles de eficiencia media más bajos que la de cantidad (TE$_I$ = 0.683 vs.\ TE$_q$ = 0.731 en el Diseño A). Existe, por tanto, un mayor margen de mejora en la dimensión de intensidad que en la de cantidad. Sin embargo, los sistemas de información hospitalaria españoles no incluyen de forma rutinaria indicadores de intensidad que permitan un seguimiento adecuado.

La solicitud de datos del RAE-CMBD con contadores de procedimientos por episodio, realizada en el marco de esta investigación, es una condición necesaria para construir medidas de intensidad más precisas. Los reguladores sanitarios deberían invertir en sistemas de información que capturen no solo cuántos pacientes se atienden, sino con qué nivel de recursos diagnósticos y terapéuticos por episodio.

\subsection{El papel de \texorpdfstring{$\psi_{12}$}{psi\_12} en la política de servicios}

La confirmación empírica de P4 ---que $\psi_{12}$ modera asimétricamente el efecto de la estructura organizativa--- tiene una consecuencia práctica directa. En los servicios quirúrgicos (donde las tareas son complementarias), la descentralización es relativamente más eficiente; en los servicios médicos no quirúrgicos (tareas sustitutas), la ventaja se reduce o desaparece. Esta heterogeneidad sugiere que las políticas de gestión hospitalaria podrían beneficiarse de un tratamiento diferenciado por tipo de servicio, permitiendo mayor autonomía en los servicios quirúrgicos y manteniendo más control en los médicos.

\section{Perspectiva internacional: lecciones de otros sistemas}
\label{sec:perspectiva_internacional}

Los dilemas que plantean los resultados de esta tesis no son exclusivos del sistema hospitalario español. La tensión entre eficiencia en volumen e intensidad de la atención, y el papel de la estructura organizativa y el mecanismo de pago en la resolución de esa tensión, son cuestiones universales que los principales sistemas sanitarios europeos han abordado con distintas soluciones institucionales.

\subsection{El NHS inglés: separación de comprador y proveedor y \textit{payment by results}}

El National Health Service (NHS) inglés introdujo en 2003--2004 el sistema de \textit{Payment by Results} (PbR), consistente en el pago prospectivo por episodio hospitalario mediante la clasificación HRG (Healthcare Resource Groups, el equivalente inglés de los GRD). Esta reforma convirtió al NHS de un sistema de financiación hospitalaria por presupuesto global en un sistema de pago por actividad, con el objetivo de incentivar la eficiencia en volumen y reducir las listas de espera.

Los resultados de la reforma PbR inglesa ilustran perfectamente el dilema predicho por nuestro modelo. \citet{street2003} y trabajos posteriores documentan que el PbR redujo los tiempos de espera y aumentó la actividad hospitalaria (evidencia de mejora en la dimensión de cantidad), pero generó preocupaciones sobre el deterioro de la calidad clínica en algunos servicios (evidencia coherente con la predicción P2). El NHS respondió introduciendo gradualmente indicadores de calidad en el sistema de pago ---el \textit{Commissioning for Quality and Innovation} (CQUIN) y las penalizaciones por reingresos--- que son exactamente los mecanismos que el modelo teórico predice como necesarios para corregir el sesgo de volumen bajo sustitución entre tareas.

\subsection{El sistema alemán de GRD (G-DRG)}

Alemania introdujo su sistema de pago hospitalario por GRD (G-DRG) en 2003, completando la transición desde el sistema de pago por día de estancia en 2004. A diferencia del NHS, el sistema alemán coexiste con una mezcla de hospitales públicos, privados con ánimo de lucro y privados sin ánimo de lucro, con una distribución aproximada de 30/35/35 en términos de número de hospitales.

\citet{herr2011} y \citet{tiemann2012} analizan la eficiencia hospitalaria alemana bajo esta diversidad de formas organizativas, encontrando que los hospitales privados con ánimo de lucro son más eficientes en costes pero no necesariamente en la calidad clínica, un resultado coherente con la asimetría P1/P2 de nuestro análisis. La experiencia alemana sugiere que la introducción del pago por GRD, en ausencia de indicadores de calidad vinculantes, puede generar los incentivos de sustitución predichos por el modelo teórico, con independencia de la forma organizativa.

\subsection{Las reformas nórdicas: \textit{patient choice} y competencia regulada}

Los países nórdicos (Suecia, Noruega, Dinamarca, Finlandia) introdujeron desde los años noventa reformas de cuasimercado hospitalario (\textit{patient choice}, competencia entre proveedores públicos, introducción de objetivos de actividad) sin privatización de la propiedad. \citet{linna1998} analiza la eficiencia técnica de hospitales finlandeses tras las reformas de los noventa, documentando mejoras de eficiencia en la dimensión de cantidad pero con evidencia menos clara en la dimensión de calidad.

La experiencia nórdica es relevante porque ilustra que es posible introducir incentivos de eficiencia en volumen dentro de una estructura organizativa predominantemente centralizada. En términos del modelo teórico, las reformas nórdicas movieron el sistema hacia esquemas de pago por actividad manteniendo la centralización de la contratación del personal médico, lo que corresponde al escenario teórico de incentivos solo sobre cantidad bajo estructura centralizada (punto de bala en la sección de implicaciones de pago fijo del Resultado~\ref{res:incentivos_centralizados}).

\subsection{Implicaciones comparadas para España}

La comparación internacional sugiere que España, con su combinación de hospitales SNS centralizados y hospitales concertados y en concesión descentralizados, representa un caso intermedio en el espectro de reformas europeas. La coexistencia de formas organizativas tan distintas en el mismo sistema, bajo financiación pública, ofrece una oportunidad de diseño institucional informado por la evidencia que otros sistemas, más homogéneos en su estructura, no tienen. Los resultados de esta tesis argumentan que esa diversidad debería gestionarse de forma consciente ---aprovechando la ventaja descentralizada en volumen en los servicios quirúrgicos y la ventaja centralizada en intensidad en los servicios médicos--- en lugar de evolucionar hacia una homogeneización de modelos en ninguna dirección.

\section{Limitaciones de las recomendaciones}
\label{sec:limitaciones_recomendaciones}

Las implicaciones de política derivadas de esta tesis deben leerse con las siguientes advertencias metodológicas.

\textbf{Primera: causalidad e identificación.} Los resultados empíricos identifican correlaciones entre estructura organizativa y eficiencia, controlando por comunidad autónoma, cluster de complejidad y tipo de pago. Sin embargo, la asignación de hospitales a estructuras descentralizadas no es aleatoria: responde a decisiones políticas regionales que pueden estar correlacionadas con características no observadas del sistema sanitario autonómico. La identificación causal requirería variación exógena en la estructura organizativa (por ejemplo, los cambios normativos que produjeron las reversiones de concesiones valencianas y madrileñas), cuyo análisis queda para trabajos futuros.

\textbf{Segunda: generalización.} La muestra incluye todos los hospitales de agudos con actividad suficiente, pero los hospitales concertados de la muestra son predominantemente de pequeño y mediano tamaño. Las recomendaciones sobre descentralización podrían no aplicarse con la misma fuerza a hospitales universitarios de tercer nivel, donde la complejidad de la cartera de servicios y la presencia de actividad docente e investigadora introducen dimensiones adicionales de output no modelizadas.

\textbf{Tercera: dinámica temporal.} El modelo teórico es estático: describe el equilibrio bajo una estructura organizativa dada. Los efectos dinámicos de un cambio de estructura (período de transición, costes de renegociación contractual, curvas de aprendizaje) no están modelizados y podrían dominar los efectos estáticos en el corto plazo.

\textbf{Cuarta: dimensiones omitidas.} El modelo se centra en dos dimensiones del output (cantidad e intensidad diagnóstica). La calidad clínica ---en el sentido de resultados de salud de los pacientes--- es una tercera dimensión crucial que no se ha podido modelizar ni medir directamente por falta de datos de seguimiento de resultados a nivel hospitalario en el SIAE. Las recomendaciones sobre estructura organizativa deberían complementarse con evidencia sobre resultados de salud que esta tesis no proporciona.

\section{Cautelas para la interpretación de las recomendaciones de política}
\label{sec:cautelas_politica}

Antes de formular recomendaciones de política, es necesario explicitar las cautelas metodológicas que deben presidir la interpretación de los resultados empíricos cuando se trasladan al plano normativo.

\subsection{Endogeneidad de la estructura organizativa y sesgo de selección}

La limitación más importante de los resultados es la posible endogeneidad de la variable $D^{desc}$. La decisión de gestionar un hospital mediante concesión administrativa, fundación pública o empresa pública no es aleatoria: responde a decisiones políticas regionales que pueden estar correlacionadas con características no observadas del sistema sanitario autonómico. En términos econométricos, si los hospitales que se descentralizan son aquellos en los que la eficiencia potencial de la descentralización es mayor (sesgo de selección positivo), el coeficiente de $D^{desc}$ en la ecuación de ineficiencia estaría sobreestimado; si, por el contrario, las regiones que descentralizan son las que tienen sistemas sanitarios más maduros y mejor gestionados (también correlacionado con mayor eficiencia), el coeficiente podría estar sesgado en cualquier dirección.

Las dummies de comunidad autónoma en la ecuación de ineficiencia absorben la heterogeneidad regional permanente, pero no la heterogeneidad específica del hospital que determina si es o no descentralizado. La identificación causal requeriría una fuente de variación exógena en la estructura organizativa, idealmente un evento normativo que cambiara la forma organizativa de algunos hospitales pero no de otros de características similares. Las reversiones de concesiones en la Comunidad Valenciana (2018) y los intentos de reversión en la Comunidad de Madrid (2015) son candidatos naturales para este tipo de análisis de diferencias en diferencias, pero en el período analizado la variación temporal inducida por estas reversiones es aún insuficiente para una identificación robusta: los hospitales revertidos son pocos y la reversión puede haberse producido precisamente porque su eficiencia era inferior a la esperada (sesgo de selección en las reversiones).

\subsection{Validez externa: ¿son generalizables los resultados?}

Los resultados de esta tesis son directamente válidos para el sistema hospitalario español en el período 2010--2023. Su generalización a otros sistemas sanitarios requiere verificar si las condiciones institucionales y tecnológicas del sistema español son comparables a las del sistema de destino. En particular:

\begin{itemize}
  \item La heterogeneidad de formas organizativas (SNS integrado, concesiones, fundaciones, concertados) es característica del modelo español; sistemas con estructuras más homogéneas (NHS inglés, Norden) tendrían menos variación en $D^{desc}$ y menor poder para identificar el efecto.
  \item El período 2010--2023 coincide con una fase de contracción del gasto sanitario público (2010--2015) y posterior recuperación. Los resultados pueden no ser representativos de períodos de expansión del gasto hospitalario.
  \item La intensidad diagnóstica medida por el SIAE ($i_{diag}$) incluye procedimientos diagnósticos pero no procedimientos terapéuticos de alta complejidad (trasplantes, cirugía cardíaca, etc.), que tienen un patrón de distribución entre hospitales muy diferente.
\end{itemize}

\subsection{Correlación versus causalidad}

La confirmación empírica de P1--P4 establece que existe una correlación estadísticamente significativa entre estructura organizativa y eficiencia hospitalaria, en la dirección predicha por el modelo teórico. Esta correlación es consistente con la hipótesis causal del modelo ---la descentralización mejora la eficiencia en cantidad porque mejora los incentivos de volumen del hospital--- pero no la establece de forma concluyente. Para establecer causalidad sería necesario identificar variación exógena en la estructura organizativa, lo que requiere diseños cuasiexperimentales que están fuera del alcance de esta tesis.

La distinción entre correlación y causalidad es especialmente importante para las recomendaciones de política: si la correlación observada no es causal, una política de descentralización que se apoyara en estos resultados podría no producir los beneficios de eficiencia esperados. La evidencia disponible es suficiente para justificar la investigación adicional con diseños causales (cuasiexperimentos), pero no para recomendar cambios de política a gran escala sin esa evidencia adicional.

\section{Hacia un diseño institucional informado por la evidencia}

En conjunto, los resultados de esta tesis apuntan a que el diseño institucional del sistema hospitalario debería tener en cuenta la multidimensionalidad del output hospitalario y la interacción entre estructura organizativa, tipo de pago y tipo de servicio. Un enfoque de ``talla única'' ---ya sea centralización total o descentralización generalizada--- ignora las complejidades que el modelo teórico pone de manifiesto y que los datos empíricos confirman.

Las reformas hospitalarias deberían evaluarse no solo por su impacto en la eficiencia global, sino por su efecto diferencial sobre las distintas dimensiones del output. Los reguladores sanitarios dispondrían de mejor información para el diseño institucional si los sistemas de evaluación del rendimiento hospitalario incorporasen medidas separadas de eficiencia en volumen y en intensidad, junto con indicadores que capturen el grado de complementariedad o sustitución entre las actividades del hospital.

Una agenda de reforma informada por los resultados de esta tesis podría incluir los siguientes elementos: (i) extensión del sistema de pago por actividad basado en GRD a toda la red hospitalaria pública y concertada, con una tarifa que diferencie explícitamente entre servicios quirúrgicos y médicos; (ii) podría resultar eficaz vincular una fracción de la financiación (por ejemplo, del orden del 5-10\%) a indicadores de intensidad diagnóstica en los contratos-programa autonómicos, para contrarrestar el efecto de sustitución predicho por el modelo; (iii) evaluación periódica multidimensional de la eficiencia hospitalaria que use técnicas de frontera con outputs separados, similar al diseño de esta tesis, como herramienta de regulación por comparación (\textit{yardstick competition}); (iv) en el caso de concesiones administrativas y hospitales concertados, los pliegos podrían beneficiarse de incluir cláusulas sobre estándares mínimos de intensidad diagnóstica y auditoría periódica de su cumplimiento.



% ###########################################################
% CAPÍTULO 8: CONCLUSIONES
% ###########################################################
